% =====================================================================
% 別冊『Git & GitHub セットアップ — はじめる前の環境準備』
% エンジン: LuaLaTeX（jlreq / book）
% ビルド:   latexmk -lualatex main_setup.tex
%
% 対象: これから本編『はじめてのバージョン管理』を読む人。
% Gitのインストール・初期設定・GitHubアカウント・gh認証だけを扱う。
% =====================================================================
\documentclass[book, paper=a4, fontsize=10.5pt, titlepage]{jlreq}

% =====================================================================
% preamble.tex — パッケージ読込・色・カスタム環境・マクロの定義
% 本文（chapters/*.tex）から使う見た目の部品はすべてここに集約する。
% エンジン: LuaLaTeX + jlreq
% =====================================================================

% --- 基本パッケージ ---------------------------------------------------
\usepackage{graphicx}          % 画像の取り込み
\usepackage{xcolor}            % 色指定
\usepackage{booktabs}          % 罫線のきれいな表
\renewcommand{\arraystretch}{1.25} % 表の行を少しゆったりさせる
\usepackage{tabularx}          % 幅を自動調整する表
\usepackage{enumitem}          % 箇条書きの調整
\usepackage{fancyvrb}          % コード枠に使う Verbatim
\usepackage[most]{tcolorbox}   % 装飾枠（コード・注意書き）
\usepackage{adjustbox}         % 図が本文幅を超えたときの自動縮小
\usepackage{tikz}              % 概念図・フロー図
\usetikzlibrary{arrows.meta,positioning,shapes.geometric,calc,fit,backgrounds}
\usepackage{float}             % 図を本文の位置に固定する [H]
\usepackage{caption}           % 図番号つきキャプションの体裁
\captionsetup{font=small, labelfont=bf, labelsep=quad}

% 画像の探索パス（既存の Markdown 版と同じ images/ を共有する）
\graphicspath{{../images/}{./images/}}

% --- 配色 -------------------------------------------------------------
% Markdown 版の Mermaid 図の色をそのまま踏襲する（薄い色は使わない方針）。
\definecolor{gitRed}{HTML}{C0392B}     % 作業ツリー・注意
\definecolor{gitOrange}{HTML}{D68910}  % ステージング
\definecolor{gitGreen}{HTML}{2E7D32}   % リポジトリ・完了
\definecolor{gitBlue}{HTML}{4A90D9}    % コミット・情報
\definecolor{gitPurple}{HTML}{6C3483}  % 自分側
\definecolor{gitTeal}{HTML}{117A65}    % 両方採用
\definecolor{gitAmber}{HTML}{8A6D0B}   % 判断ノード
\definecolor{codeBg}{HTML}{F4F5F7}     % コード枠の背景
\definecolor{codeFrame}{HTML}{D0D3D9}  % コード枠の枠線

% --- コード枠 ---------------------------------------------------------
% 使い方: \begin{codebox}[bash] ... Verbatim ... \end{codebox}
% 省略時のラベルは bash。右上に言語ラベルを小さく表示する。
\newtcolorbox{codebox}[1][bash]{%
  enhanced, breakable,
  colback=codeBg, colframe=codeFrame, boxrule=0.6pt, arc=2pt,
  left=9pt, right=9pt, top=13pt, bottom=8pt,
  before skip=8pt, after skip=8pt,
  overlay={%
    \node[anchor=north east, font=\sffamily\scriptsize\bfseries,
          text=gray!60] at ([xshift=-6pt, yshift=-3pt]frame.north east) {#1};%
  },
}
% コード本体を書くための Verbatim 環境（小さめ・行間やや詰め）
\DefineVerbatimEnvironment{code}{Verbatim}{fontsize=\small, xleftmargin=0pt}

% 本文中のインラインコード: \code{git status}
\newcommand{\cmd}[1]{\texttt{#1}}

% --- 注意書き（admonition）枠 ----------------------------------------
% ヒント（緑）・注意（赤）・補足（グレー）の3種。emoji は使わずラベルで表す。
% タイトルは本文（明るい背景）に直接乗るため、白文字だと見えない。
% タイトル文字色は各テーマ色にして、はっきり読めるようにする。
\newtcolorbox{hintbox}{%
  enhanced, breakable,
  colback=gitGreen!5, colframe=gitGreen, boxrule=0.8pt, arc=2pt,
  left=8pt, right=8pt, top=6pt, bottom=6pt,
  coltitle=gitGreen, fonttitle=\sffamily\bfseries\small,
  title=ヒント, attach title to upper=\quad,
}
\newtcolorbox{warnbox}{%
  enhanced, breakable,
  colback=gitRed!5, colframe=gitRed, boxrule=0.8pt, arc=2pt,
  left=8pt, right=8pt, top=6pt, bottom=6pt,
  coltitle=gitRed, fonttitle=\sffamily\bfseries\small,
  title=注意, attach title to upper=\quad,
}
% 執筆メモ（雛形用の枠）。原稿を清書したら、この枠ごと削除する。
\newtcolorbox{draftbox}{%
  enhanced, breakable,
  colback=gitAmber!6, colframe=white, boxrule=0pt, arc=2pt,
  borderline={0.9pt}{0pt}{gitAmber, dashed},
  left=8pt, right=8pt, top=6pt, bottom=6pt,
  coltitle=gitAmber, fonttitle=\sffamily\bfseries\small,
  title=執筆メモ, attach title to upper=\quad,
}
\newtcolorbox{notebox}{%
  enhanced, breakable,
  colback=gray!8, colframe=gray!60!black, boxrule=0.8pt, arc=2pt,
  left=8pt, right=8pt, top=6pt, bottom=6pt,
  coltitle=gray!30!black, fonttitle=\sffamily\bfseries\small,
  title=補足, attach title to upper=\quad,
}

% --- スクリーンショット取り込み --------------------------------------
% \screenshot[幅の割合]{ファイル}{キャプション}
% フロートにせず本文の流れに沿って中央寄せで置く（説明の直後に見せたいため）。
\newcommand{\screenshot}[3][0.7]{%
  \begin{center}
    \includegraphics[width=#1\linewidth]{#2}\\[3pt]
    {\small\itshape #3}%
  \end{center}%
}

% \screenshottrim[幅の割合]{ファイル}{キャプション}{trim指定}
% 余白の多いスクリーンショットを切り詰めて大きく見せる。
% trim は「左 下 右 上」の順で、画像1ピクセル = 1bp として指定する
% （例: 1920x1080 の下半分を捨てる → {0 540bp 0 0}）。
\newcommand{\screenshottrim}[4][0.95]{%
  \begin{center}
    \includegraphics[width=#1\linewidth, trim=#4, clip]{#2}\\[3pt]
    {\small\itshape #3}%
  \end{center}%
}

% --- TikZ 共通スタイル ------------------------------------------------
% Mermaid のフロー図を再現するためのノード／矢印スタイル。
\tikzset{
  flownode/.style={%
    rounded corners=4pt, draw=none, text=white,
    font=\sffamily\small, align=center,
    minimum height=1.15cm, minimum width=2.2cm, inner sep=6pt},
  decisionnode/.style={%
    diamond, aspect=1.9, draw=none, text=white,
    font=\sffamily\small, align=center, inner sep=3pt},
  flowarrow/.style={-{Latex[length=3.6mm, width=3mm]}, line width=1.8pt, draw=black},
  edgelabel/.style={font=\sffamily\normalsize\bfseries, fill=white, inner sep=3pt},
}
% 図全体を中央寄せで置くためのヘルパ環境（図番号なし）。
% 本文幅を超えた図は自動で縮小する（幅が足りている図は拡大しない）。
\newenvironment{diagram}%
  {\par\medskip\begin{center}\begin{adjustbox}{max width=\linewidth}}%
  {\end{adjustbox}\end{center}\par\medskip}

% --- 図番号つきの図（本文が「図1.2のように」と参照する概念図・フロー図用）---
% 手順に紐づく操作スクショには使わない（そちらは \screenshot を使う）。
\makeatletter
% TikZ など本体を囲む環境: \begin{numfig}{キャプション}{ラベル} ... \end{numfig}
\newenvironment{numfig}[2]{%
  \def\nf@cap{#1}\def\nf@lab{#2}%
  \begin{figure}[H]\centering\begin{adjustbox}{max width=\linewidth}%
}{%
  \end{adjustbox}\caption{\nf@cap}\label{\nf@lab}\end{figure}%
}
\makeatother
% 画像1枚を図番号つきで置く: \figimage[幅]{ファイル}{キャプション}{ラベル}
\newcommand{\figimage}[4][0.95]{%
  \begin{figure}[H]\centering
    \includegraphics[width=#1\linewidth]{#2}
    \caption{#3}\label{#4}
  \end{figure}%
}


% hyperref はプリアンブルの最後（他パッケージ読込後）に置くのが安全。
\usepackage[unicode, hidelinks]{hyperref}

% --- 書誌情報 ---------------------------------------------------------
\title{Git \& GitHub セットアップ}
\author{}
\date{\today}

\begin{document}

% --- 表紙 -------------------------------------------------------------
\begin{titlepage}
  \centering
  \vspace*{4cm}
  {\sffamily\bfseries\Huge Git \& GitHub セットアップ\par}
  \vspace{1.2cm}
  {\sffamily\Large はじめる前の環境準備\par}
  \vspace{0.6cm}
  {\sffamily\normalsize ── 本編『はじめてのバージョン管理』を始める前に ──\par}
  \vfill
  {\large \today\par}
\end{titlepage}

% --- この本のねらい ---------------------------------------------------
\thispagestyle{empty}
\vspace*{1.5cm}
{\sffamily\LARGE\bfseries はじめに —この冊子でやること—\par}
\vspace{1em}

本書は、Git・GitHub を使いはじめる前の\textbf{環境の準備だけ}をまとめた冊子です。\textbf{一度やれば終わり}の作業なので、済んだらこの冊子は棚にしまって構いません。

\begin{itemize}[leftmargin=1.6em]
  \item 自分のPCに Git をインストールする
  \item コミットに記録される名前とメールアドレスを登録する
  \item GitHub のアカウントを作る
  \item \cmd{gh}（GitHub CLI）をインストールし、GitHub にログインする
\end{itemize}

\begin{hintbox}
\textbf{Gitの使い方そのもの}（コミットの取り方、GitHubへのアップロード）は、本編『はじめてのバージョン管理』で扱います。準備が終わったらそちらへ進んでください。
\end{hintbox}

\clearpage

% --- 目次 -------------------------------------------------------------
\tableofcontents
\clearpage

% --- 本文 -------------------------------------------------------------
% =====================================================================
% セットアップ編 第1章 Gitを使えるようにする
% 元原稿: chapters/01_git.tex（本編から分離）
% =====================================================================
\chapter{Gitを使えるようにする}

この章では、\textbf{自分のPCでGitを使えるようにするまで}を扱います。やることは「Gitをインストールする」と「名前とメールアドレスを登録する」の2つだけです。設定はPC全体に保存されるので、\textbf{ここは一度やれば終わり}です。この章を終えたら、本編『はじめてのバージョン管理』に進んでください。


% ---------------------------------------------------------------------
\section{Gitのインストール}

まず、Gitが\textbf{すでに入っていないか}を確認しましょう。ターミナル（WindowsはPowerShellやコマンドプロンプト）で次を実行します。

\begin{codebox}[bash]
\begin{code}
git --version
# 出力例: git version 2.45.0
\end{code}
\end{codebox}

\textbf{この結果によって、次に進む先が変わります。}

\begin{itemize}[leftmargin=1.4em]
  \item \textbf{バージョンが表示された人}（例：\cmd{git version 2.45.0}）\\
        → Gitは\textbf{すでに入っています}。この後のOS別の手順は\textbf{すべて飛ばして}、\ref{sec:first-config}節\textbf{「最初の設定」}へ進んでください。
  \item \textbf{\cmd{command not found}（見つからない）と出た人}\\
        → Gitが\textbf{入っていません}。お使いのOSの手順へ進んでください。\\
        \textbf{Windows} → \ref{sec:install-win}項　／　\textbf{macOS} → \ref{sec:install-mac}項　／　\textbf{Linux} → \ref{sec:install-linux}項
\end{itemize}

インストールが終わったら、もう一度 \cmd{git -{}-version} を実行してください。バージョンが表示されれば成功です。

\begin{warnbox}
\textbf{インストール直後は、開いていたターミナルではまだ \cmd{git --version} が反映されないことがあります。} PATH（コマンドの置き場所を示す設定）がターミナルの起動時に読み込まれるためです。\cmd{command not found} のままなら、\textbf{ターミナルをいったん閉じて開き直して}から、もう一度実行してください。
\end{warnbox}

\subsection{Windows}
\label{sec:install-win}

Windowsでは「Git for Windows」をインストールします。これには \textbf{Git Bash} という、macOS/Linux風のコマンドが使えるターミナルが付いてきます。本シリーズのWindowsの説明は、この Git Bash を使う前提です。

\begin{enumerate}[leftmargin=1.6em]
  \item \url{https://git-scm.com/download/win} からインストーラをダウンロードし、ダブルクリックで起動します。
  \item インストーラの画面が続きます。\textbf{基本はすべてデフォルトのまま「Next」で進めて問題ありません。} ただし次の3つの画面だけは選択を確認してください。
\end{enumerate}

\noindent\textbf{① 既定のエディタ}

デフォルトはVimですが不慣れだと難しいため、VS Codeを使う予定なら \textbf{「Use Visual Studio Code as Git's default editor」} を選びます。
\screenshot[0.62]{01_git/win_install_05.png}{既定エディタの選択}

\noindent\textbf{② 初期ブランチ名}

\textbf{「Override the default branch name…」} を選び、\cmd{main} と入力します（世界的に \cmd{main} が標準）。
\screenshot[0.62]{01_git/win_install_06.png}{初期ブランチ名の設定}

\noindent\textbf{③ 改行コードの扱い}

WindowsとほかのOSで改行コードが異なるため重要です。複数環境で開発するなら \textbf{「Checkout as-is, commit Unix-style line endings」} を選びます。
\screenshot[0.62]{01_git/win_install_10.png}{改行コードの設定}

残りの画面はデフォルトのまま進め、最後に「Install」をクリックします。インストール後、スタートメニューから \textbf{Git Bash} を起動し、\cmd{git --version} でバージョンが出れば完了です。

\subsection{macOS}
\label{sec:install-mac}

次のどちらか一方でOKです。

\noindent\textbf{方法A: Xcode Command Line Tools（最も手軽）}

ターミナル（\cmd{アプリケーション > ユーティリティ > ターミナル.app}）を開き、次を実行します。

\begin{codebox}[bash]
\begin{code}
xcode-select --install
\end{code}
\end{codebox}

ダイアログが出たら「インストール」をクリックします。

\noindent\textbf{方法B: Homebrew（最新版を使いたい場合）}

Homebrew（\url{https://brew.sh/index_ja}）が入っていれば、次で最新のGitが入ります。

\begin{codebox}[bash]
\begin{code}
brew install git
\end{code}
\end{codebox}

どちらの場合も、\cmd{git --version} でバージョンが表示されれば完了です。

\subsection{Linux}
\label{sec:install-linux}

お使いのディストリビューションのパッケージマネージャを使います。自分の環境に合うものを1つ選んで実行します。

\begin{codebox}[bash]
\begin{code}
# Debian / Ubuntu 系
sudo apt update && sudo apt install git

# Fedora / RHEL 系
sudo dnf install git

# Arch 系
sudo pacman -S git
\end{code}
\end{codebox}

\cmd{git --version} でバージョンが表示されれば完了です。

% ---------------------------------------------------------------------
\section{最初の設定}
\label{sec:first-config}

インストール直後に\textbf{一度だけ}行う設定です。コミットに記録される「あなたの名前とメールアドレス」を登録します。ここで設定した名前・メールが、後でGitHub上の履歴にも表示されます。

ターミナル（Windowsは Git Bash）を開き、次を順に実行します。

\begin{codebox}[bash]
\begin{code}
# あなたの名前を登録（GitHubのアカウント名に合わせるのがおすすめ）
git config --global user.name "your-name"

# あなたのメールアドレスを登録（GitHubで使う予定のメール）
git config --global user.email "you@example.com"

# 設定内容を確認する
git config --global --list
\end{code}
\end{codebox}

\begin{warnbox}
名前やメールにスペースが含まれる場合は、上のように \textbf{ダブルクォート \cmd{"} で囲んでください}。
\end{warnbox}

\cmd{--global} は「このPC全体の共通設定」という意味です。一度やれば、以降すべてのリポジトリに適用されます。

\begin{hintbox}
おすすめの追加設定（新規リポジトリの初期ブランチ名を \cmd{main} に統一）
\begin{codebox}[bash]
\begin{code}
git config --global init.defaultBranch main
\end{code}
\end{codebox}
\end{hintbox}


% =====================================================================
% セットアップ編 第2章 GitHubを使えるようにする
% 元原稿: chapters/02_github.tex（本編から分離）
% =====================================================================
\chapter{GitHubを使えるようにする}

この章では、\textbf{GitHubのアカウントを作り、自分のPCからGitHubにログインするまで}を扱います。ここまで済ませておくと、本編でGitHubに\cmd{push}（アップロード）するとき、認証で立ち止まらずに済みます。前章と同じく、\textbf{ここも一度やれば終わり}です。

なお、GitHubがどういうサービスなのかという説明は、本編『はじめてのバージョン管理』の第2章で行います。この章は\textbf{準備だけ}に集中します。

% ---------------------------------------------------------------------
\section{GitHubアカウントを作る}

まだアカウントがなければ作成します。

\begin{enumerate}[leftmargin=1.6em]
  \item \url{https://github.com/} を開き、\textbf{「Sign up」} をクリックします。
  \item メールアドレス・パスワード・ユーザー名を入力して進めます。
  \item 登録したメールアドレスに届く確認コードを入力すれば完了です。
\end{enumerate}

\begin{hintbox}
ここで登録する\textbf{メールアドレスは、前章の \cmd{git config -{}-global user.email} で設定したものに揃える}のがおすすめです。揃えておくと、コミットとGitHubアカウントが正しく結びつきます。
\end{hintbox}

% ---------------------------------------------------------------------
\section{GitHubにログインする（認証）}

GitHubに \cmd{push}（アップロード）するときは、「本当にあなたか」を確認する\textbf{認証}が必要です。本シリーズでは、いちばん手軽な \textbf{GitHub CLI（\cmd{gh}）でのログイン}を使います。これで一度ログインしておくと、\cmd{gh} がGitの認証を裏で肩代わりしてくれるため、\textbf{トークンを自分で発行しなくても \cmd{push} / \cmd{pull} がそのまま通ります}。

\begin{hintbox}
SSH鍵やトークン（PAT）を手動で使う方式もあります。頻繁に使うようになったら、『逆引きチートシート』の「認証の別方式（SSH / PAT）」を参照してください。まずは \cmd{gh} だけでOKです。
\end{hintbox}

\noindent\textbf{手順1：\cmd{gh} をインストールする}

\begin{warnbox}
\textbf{\cmd{gh} はGit本体には含まれていません。} 第1章でGitを入れても \cmd{gh} は入らないので、別途インストールが必要です。まず \cmd{gh -{}-version} を実行し、\cmd{command not found} なら下記で入れてください。
\end{warnbox}

\begin{codebox}[bash]
\begin{code}
# macOS（Homebrew）
brew install gh

# Windows（winget）
winget install --id GitHub.cli

# Linux（Debian / Ubuntu 系）
sudo apt install gh
\end{code}
\end{codebox}

インストール後、\cmd{gh --version} でバージョンが表示されれば成功です（うまくいかない場合は公式手順 \url{https://cli.github.com/} を参照）。

\noindent\textbf{手順2：ログインする}

\begin{codebox}[bash]
\begin{code}
gh auth login
\end{code}
\end{codebox}

対話形式で聞かれるので、次のように選びます。

\begin{center}
\begin{tabularx}{\linewidth}{@{}l X@{}}
\toprule
質問 & 選ぶもの \\
\midrule
Where do you use GitHub? & \textbf{GitHub.com} \\
preferred protocol for Git operations? & \textbf{HTTPS} \\
Authenticate Git with your GitHub credentials? & \textbf{Y} \\
How would you like to authenticate? & \textbf{Login with a web browser} \\
\bottomrule
\end{tabularx}
\end{center}

質問文や順序は \cmd{gh} のバージョンによって少し変わることがあります。上の表と完全に一致しなくても、\textbf{同じ意味の選択肢}（GitHub.com／HTTPS／ブラウザでログイン）を選べば問題ありません。

最後の選択でワンタイムコードが表示され、Enterを押すとブラウザが開きます。表示されたコードを入力し、\textbf{Continue} を押せば認証完了です。

\screenshot[0.62]{02_github/gh_auth_login.png}{GitHubログイン画面（Continueを押す）}
\screenshot[0.62]{02_github/gh_authorize_window.png}{GitHub認証画面（ワンタイムコードを入力）}

これで \cmd{gh} がGitの認証を肩代わりし、\textbf{以降 \cmd{git push} / \cmd{git pull} で認証を聞かれなくなります}。

\begin{codebox}[bash]
\begin{code}
# ログイン状態を確認する
gh auth status
\end{code}
\end{codebox}

\begin{warnbox}
トークンの有効期限が切れたり、GitHub側でアクセスを取り消したりすると、認証を聞かれるようになります。そのときは \cmd{gh auth login} で入り直してください。
\end{warnbox}


% ---------------------------------------------------------------------
\section{セットアップ完了チェック}

ここまでで準備は終わりです。次の4つを実行して、\textbf{すべてエラーなく表示されれば完了}です。本編『はじめてのバージョン管理』へ進んでください。

\begin{codebox}[bash]
\begin{code}
git --version              # Gitのバージョンが表示される
git config --global --list # user.name と user.email が表示される
gh --version               # ghのバージョンが表示される
gh auth status             # Logged in to github.com と表示される
\end{code}
\end{codebox}

\begin{center}
\begin{tabularx}{\linewidth}{@{}l X@{}}
\toprule
確認すること & うまくいかないときは \\
\midrule
Gitが入っているか & 第1章「Gitのインストール」へ戻る \\
名前・メールが登録されているか & 第1章「最初の設定」へ戻る \\
\cmd{gh} が入っているか & 本章「GitHubにログインする」の手順1へ戻る \\
GitHubにログインできているか & \cmd{gh auth login} をもう一度実行する \\
\bottomrule
\end{tabularx}
\end{center}

インストールしたはずなのに \cmd{command not found} と出るときは、\textbf{ターミナルをいったん閉じて開き直して}から、もう一度試してください。インストール直後は、設定がターミナルに反映されていないことがあります。


\end{document}
